\documentclass[UTF8]{ctexart}

% 支持从项目根目录、强化学习目录或当前文档目录读取共享样式。
\makeatletter
\def\input@path{{../}{../../}}
\makeatother
\usepackage{researchnotes}

\title{多臂赌博机与强化学习}
\author{}
\date{}

\begin{document}
\maketitle

\section{共同目标：通过反馈学习决策}
\label{sec:bandit-rl-objective}

\keyterm{强化学习}（reinforcement learning）研究如何通过行动获得反馈，学习使累计奖励更高的决策方式。
\keyterm{多臂赌博机}（multi-armed bandit）保留了这一机制，并简化了环境状态变化与延迟奖励问题。
以下讨论经典的平稳随机多臂赌博机。

设有 $K$ 根拉杆，第 $t$ 轮选择 $A_t\in\{1,\ldots,K\}$，随后观察随机奖励 $R_t$。
给定本轮选择后，奖励分布只由所选拉杆决定，不随时间或过去记录变化，且期望有限；这些分布事先未知。
\keyterm{策略}（policy）$\pi$ 根据此前的选择与奖励记录决定下一次行动，目标是在给定的 $T$ 轮中实现
\begin{equation}
  \max_{\pi}\;\mathbb{E}_{\pi}\!\left[\sum_{t=1}^{T}R_t\right].
  \label{eq:bandit-rl-return}
\end{equation}
期望同时考虑奖励的随机性与策略可能采用的随机选择。

\section{简化之处：没有动作引起的状态变化}
\label{sec:bandit-rl-simplification}

在一般强化学习中，动作既影响当前奖励，也可能改变下一状态，从而影响后续收益。
例如，机器人为了绕过障碍，可能暂时远离出口；当前没有收益的动作，也可能为最终成功创造条件。
\keyterm{时间信用分配}（temporal credit assignment）就是判断过去哪些动作应当为后来的结果负责。

多臂赌博机每轮结束后仍面对相同的拉杆及其奖励分布，本轮奖励直接对应本轮选择，
无需沿环境状态序列追溯延迟奖励。因此，它可建模为\keyterm{只有一个状态的马尔可夫决策过程}
（Markov decision process，MDP）：拉杆对应动作，每次行动后都回到同一状态。

\section{保留的难点：探索与利用}
\label{sec:bandit-rl-exploration}

\keyterm{利用}（exploitation）是选择目前估计收益最高的拉杆；
\keyterm{探索}（exploration）是尝试了解较少的拉杆，以判断它是否更好。
探索可能牺牲当前奖励，却改善后续决策，因此只追求眼前估计收益最高，并不一定使式\eqref{eq:bandit-rl-return}最大。

\keyterm{行动不改变拉杆，却会改变决策者对拉杆的认识}。
所以，一个环境状态并不意味着各轮决策彼此独立：历史反馈仍会影响下一次选择。
多臂赌博机的简单之处在于省去了复杂的状态转移，同时保留了通过试错学习以及权衡探索与利用的核心问题。

\end{document}
